\section{Supplementary Problem 2.1}
The reference problem defines the discrete-time signal
\begin{equation}
  x(n)=\begin{cases}
    1+\dfrac{n}{3}, & -3\leq n\leq-1,\\
    1, & 0\leq n\leq3,\\
    0, & \text{otherwise}.
  \end{cases}
\end{equation}
All time indices in this problem are integers.

\subsection{Part (a): Sample values}
Substituting the integer time indices gives
\begin{table}[H]
  \centering
  \caption{Sample values of the piecewise signal}
  \label{tab:p21-input}
  \begin{tabular}{lrrrrrrr}
    \toprule
    $n$ & $-3$ & $-2$ & $-1$ & 0 & 1 & 2 & 3\\
    \midrule
    $x(n)$ & 0 & $1/3$ & $2/3$ & 1 & 1 & 1 & 1\\
    \bottomrule
  \end{tabular}
\end{table}
The sample at $n=-3$ is zero, since $1+(-3)/3=0$.
One vector representation is
\[
  \texttt{xn}=[0,1/3,2/3,1,1,1,1],\qquad \texttt{xorigin}=4.
\]

\subsection{Part (b): Changing the operation order}
\subsubsection{Fold first, then delay by four samples}
Let $v(n)=x(-n)$. Delaying this intermediate signal gives
\begin{equation}
  y_{b1}(n)=v(n-4)=x(-(n-4))=x(4-n).
\end{equation}
Equivalently, an original sample at time $m$ moves first to $-m$ and
then to $4-m$. The resulting piecewise expression is
\begin{equation}
  y_{b1}(n)=\begin{cases}
    1, & 1\leq n\leq4,\\
    (7-n)/3, & 5\leq n\leq7,\\
    0, & \text{otherwise}.
  \end{cases}
\end{equation}

\subsubsection{Delay first, then fold}
Let $w(n)=x(n-4)$. Folding this intermediate signal gives
\begin{equation}
  y_{b2}(n)=w(-n)=x(-n-4).
\end{equation}
An original sample at time $m$ moves first to $m+4$ and then to
$-m-4$. Therefore,
\begin{equation}
  y_{b2}(n)=\begin{cases}
    1, & -7\leq n\leq-4,\\
    (-n-1)/3, & -3\leq n\leq-1,\\
    0, & \text{otherwise}.
  \end{cases}
\end{equation}
\begin{table}[H]
  \centering
  \caption{Transformed samples for the two operation orders}
  \label{tab:p21-transforms}
  \begin{tabular}{lrrrrrrr}
    \toprule
    Time for $y_{b1}$ & 1 & 2 & 3 & 4 & 5 & 6 & 7\\
    Time for $y_{b2}$ & $-7$ & $-6$ & $-5$ & $-4$ & $-3$ & $-2$ & $-1$\\
    \midrule
    Corresponding value & 1 & 1 & 1 & 1 & $2/3$ & $1/3$ & 0\\
    \bottomrule
  \end{tabular}
\end{table}

\subsection{Part (c): The signal x(-n+4)}
Since $-n+4=4-n$, the required signal is exactly $y_{b1}(n)$.
Its plot and values coincide with the fold-then-delay result in part (b).

\subsection{Part (d): General rule}
For $y(n)=x(-n+k)=x(-(n-k))$, fold the original signal and then shift
the folded signal to the right by $k$ samples when $k>0$.
Equivalently, advance the original signal by $k$ samples and then fold.
The direct sample mapping is
\begin{equation}
  m\longmapsto k-m.
\end{equation}
For negative $k$, the shift directions reverse. Folding and a nonzero
time shift do not commute in general; changing their order reverses
the direction of the resulting shift. Here the original sample at $m=0$
lands at $n=4$ in $y_{b1}$ but at $n=-4$ in $y_{b2}$.

\subsection{Part (e): Impulse and unit-step representations}
Define $\delta(n)=1$ for $n=0$ and zero otherwise, and $u(n)=1$ for
$n\geq0$ and zero otherwise. Using a weighted impulse at each nonzero
sample gives
\begin{align}
  x(n)={}&\tfrac13\delta(n+2)+\tfrac23\delta(n+1)
          +\delta(n)+\delta(n-1)\notag\\
       &+\delta(n-2)+\delta(n-3).
\end{align}
The two piecewise intervals can also be selected by differences of
unit steps:
\begin{equation}
  x(n)=\left(1+\frac n3\right)\bigl[u(n+3)-u(n)\bigr]
       +u(n)-u(n-4).
\end{equation}
The first step difference selects $-3\leq n\leq-1$, and the second
selects $0\leq n\leq3$, reproducing the original definition.

\subsection{Scilab source and expected checks}
The script evaluates the original piecewise function at each transformed
time argument. It plots parts (a), (b1), (b2), and (c), and compares the
two reconstructions from part (e) against the original samples on
$-8\leq n\leq8$. This interval includes all nonzero samples of all
plotted signals. Parts (b1) and (c) should agree exactly; reconstruction
errors should be zero up to floating-point rounding.
\lstinputlisting[caption={Scilab implementation of supplementary Problem 2.1},label={lst:p21}]{../2_1.sci}

\subsection{Execution evidence}
\resultfigure{images/problem21.png}
  {Original signal and the three transformations in Problem 2.1}
  {fig:p21}
The captured plots agree with the calculated sample values. In particular,
parts (b1) and (c) coincide on the positive-time side, while part (b2)
appears at the corresponding negative time indices. The console
reconstruction errors are separate checks and are not shown in this figure.
